% META: {"id": "1NSI-PM-CO-001", "chapitre": "1NSI-PROJET-METHODES", "type_objet": "corrige", "exercice_ref": "1NSI-PM-EX-001", "capacites": ["P-HIST-01"], "sources_inspiration": [], "status": "verified"}
% BEGIN-VERIFY
% evenements = [
%     (1642, "Pascaline, machine a calculer mecanique", "Blaise Pascal"),
%     (1936, "Concept de machine universelle", "Alan Turing"),
%     (1971, "Premier microprocesseur commercial", "Intel"),
%     (1991, "Premiere version publique de Python", "Guido van Rossum"),
% ]
% assert evenements == sorted(evenements, key=lambda e: e[0])
% def evenements_entre(evenements, debut, fin):
%     return [e for e in evenements if debut <= e[0] <= fin]
% resultat = evenements_entre(evenements, 1900, 1980)
% assert [e[0] for e in resultat] == [1936, 1971]
% END-VERIFY
\begin{corrige}{1NSI-PM-EX-001}
\textbf{1.} \lstinline{evenements == sorted(evenements, key=lambda e: e[0])} vaut
\lstinline{True} : la liste est bien triée par année croissante.

\textbf{2.} Entre 1900 et 1980 : le concept de machine universelle de Turing (1936) et le
premier microprocesseur commercial (1971).

\textbf{3.} La Pascaline (1642) est une machine \textbf{mécanique} conçue pour une tâche
unique (l'addition) : elle ne peut exécuter qu'un seul type de calcul, câblé dans sa
mécanique. Le concept de machine universelle de Turing (1936) est une avancée
\textbf{théorique} majeure : il montre qu'une seule machine peut, en principe, exécuter
\textbf{n'importe quel} algorithme, à condition de lui fournir le bon programme — c'est le
principe fondateur de l'ordinateur moderne, programmable.
\end{corrige}
